% Detection, Forecasting, and Clinical Readiness
\subsection{Detection, Forecasting, and Clinical Readiness}
\label{sec:detection-forecasting-readiness}

Detection and forecasting serve different clinical purposes. Detection aims to identify seizures as they occur for timely intervention, while forecasting aims to predict seizures before onset to allow preventive action. These different objectives lead to distinct technical approaches, validation requirements, and clinical readiness levels.

\subsubsection{Maturity Comparison}

Detection shows greater clinical maturity than forecasting. Two detection studies meet Phase 3 FPR benchmarks defined by the ILAE, though both are Phase 2 studies themselves. \textcite{spahrDeepLearningbasedDetection2025} achieved 0.0054/h FPR (approximately 1 false alarm per 8 days) with 96\% sensitivity in a prospective multi-center study of 384 patients. In a commentary on Larsen et al. 2024, \textcite{fineDetectionKeyAutomated2025} described 0.023/h FPR with 100\% sensitivity on an independent test set of 3 patients with 10 tonic seizures. Commercial devices exist for detection including the Embrace watch with FDA 510(k) clearance and NightWatch with CE approval in Europe.

Forecasting shows consistent AUC around 0.74 to 0.75 across studies, suggesting a performance ceiling for non-invasive approaches. \textcite{meiselMachineLearningWristband2020} reported AUC 0.74 with IoC of 14.1\% over chance. \textcite{nasseriAmbulatorySeizureForecasting2021} reported AUC 0.80 (mean) but studied only 6 patients. \textcite{vielufSeizureMonitoringCombined2025} achieved AUC 0.75 in a larger sample of 70 patients. However, no forecasting study meets a clear clinical benchmark and no device has regulatory approval for seizure forecasting.

Detection benefits from clearer clinical requirements. The ILAE has defined guidelines for detection device validation including FPR targets below 0.05 to 0.1 per hour. Forecasting lacks equivalent standardized guidelines. This gap complicates assessment of forecasting readiness and clinical utility.

Sample size disparities further reflect the maturity difference. Detection studies have a median sample size of 66 participants compared to 40 for forecasting studies. Larger detection studies enable more robust performance estimates. The largest detection study enrolled 384 patients across eight centers, while the largest forecasting study enrolled 70 patients.
